%% kapa: exam=1 round=22 year=2011 period=1 form=A subject=economics subject_ko=경제원론 seq=34 number=74
%% source: Q-Net 감정평가사 기출문제 (cst00309) artlSeq=5036790 file=extracted/1008/5036790/2011년 제22회 감정평가사 1차시험 기출문제/1교시.hwp
%% conversion: kordoc 4.24.51 -> tools/gen1.py (auto); review=proofread-sonnet; flags=text-verified-hwp
\begin{question}{74}
고전학파의 대부자금설이 성립할 경우 정부가 저축을 촉진하기 위해 이자소득세를 인하하고 동시에 투자를 촉진하는 투자세액공제제도를 도입할 때 예상되는 대부자금 시장의 변화로 옳은 것은? (단, 수요곡선은 우하향, 공급곡선은 우상향)
\begin{choices}
  \item 균형이자율 상승, 균형거래량 증가
  \item 균형이자율 상승, 균형거래량 감소
  \item 균형이자율 하락, 균형거래량 증가
  \item 균형이자율 하락, 균형거래량 증감 불분명
  \item 균형이자율 등락 불분명, 균형거래량 증가
\end{choices}
\end{question}
\answer{5}
